\documentclass[11pt]{article}

% parameters are fairly self-explanatory here, thankfully
\usepackage[width=7in,height=9.5in,top=0.75in,centering]{geometry}

% for math-ing
\usepackage{amsmath}
\usepackage{amssymb}


%%%%%% tables %%%%%%%%

\usepackage{array}
\renewcommand{\arraystretch}{1.5}
\setlength{\tabcolsep}{8pt}

%%%%%% plotting %%%%%%

\usepackage[dvipsnames]{xcolor}
\usepackage{tikz}
\usepackage{pgfplots}

\pgfplotsset{every axis/.append style={
                    axis x line=middle,    % put the x axis in the middle
                    axis y line=middle,    % put the y axis in the middle
                    axis line style={<->,color=black}, % arrows on the axis
                    xlabel={$x$},          % default put x on x-axis
                    ylabel={$y$},          % default put y on y-axis
            }}
            
 \pgfplotsset{compat=1.14} % sometimes, everything falls apart without this
 
 %\usepgfplotslibrary{fillbetween} %for shading between curves

%%%%%% commands %%%%%%

% exam heading; course name is first argument, exam information is second argument
\newcommand{\Heading}[2]{
    \fbox{
        {\Large\textbf{#1}}
        }
    \hfill
    \fbox{
        \begin{minipage}{0.2\textwidth}
            \begin{center}
            \textbf{#2}
            \end{center}
        \end{minipage}
    }
    \par\vspace{0.1in}}

% for being lazy while beautifying integrals
\newcommand{\dint}{\displaystyle\int}

% for being lazy while beautifying limits
\newcommand{\dlim}{\displaystyle\lim}

%%%%%% stuff %%%%%%

\usepackage{fancyhdr}

%\setcounter{page}{0} % first page is page 0

\parindent0pt % no indent

\fboxsep=0.25cm % little space in fbox border

%%%%%% BEGIN! %%%%%%%

\begin{document}

\pagestyle{fancy}

% record goals here to create sense of transparency

\fancyhead[l]{%
  \ifcase\value{page}
  % empty test for page = 0
  \or {\bf\color{Fuchsia} I1: Meaning of an Integral}% page = 7
  \or {\bf\color{Fuchsia} I2: Properties of Definite Integrals}% page = 8
  \else
  % Empty chead here!
  \fi
}


\Heading{MATH 1190/1210 Survey of Calculus I}{Quiz \\Learning Targets I1-I2\\Fall 2025}

\vspace{0.1in}

\textbf{Name} \hrulefill \, \begin{minipage}{0.16\textwidth}\begin{center}\textbf{UVA ID\\ (e.g. hdj4nd)} \end{center}\end{minipage}\hrulefill\\




%%%%%%%%%%%%%%%%%%%%%%%%%%%%%%%%%%%%%%%%%%%%%%%%%%%%%
%%%%%%%%%%% BEGIN PROBLEM I1
%%%%%%%%%%%%%%%%%%%%%%%%
%%%%%%%%%%%%%%%%%%%%%%%%%%%%%%%%%%%%%%%%%%%%%%%%%%%%%

The daily usage of electricity of a household (in kWh per day) on day $t$ after January 1st is represented by $W(t)$.

\begin{enumerate}
    \item Write a definite integral that represents the total electricity used by the household during the year (365 days). If the cost of electricity is $\$0.15$ per kWh, extend your expression to represent the total annual cost of electricity for this household. Include units for both expressions
    \vfill
    \item Using the provided graph, estimate the annual cost using rectangles. 
    \begin{figure}[h]
        \centering
        \includegraphics[width=0.75\linewidth]{Fall 2025/Mikhail_Quizes/graph.png}
    \end{figure}
    \vfill
\end{enumerate}


%%%%%%%%%%%%%%%%%%%%%%%%%%%%%%%%%%%%%%%%%%%%%%%%%%%%%
%%%%%%%%%%% END PROBLEM I2 %%%%%%%%%%%%%%%%%%%%%%%%%%
%%%%%%%%%%%%%%%%%%%%%%%%%%%%%%%%%%%%%%%%%%%%%%%%%%%%%
\newpage

The graph of $f(x)$ is shown below
\begin{figure}[h]
    \centering
    \includegraphics[width=0.75\linewidth]{Fall 2025/Mikhail_Quizes/i2graph.png}
\end{figure}

Given $\displaystyle\int_{5}^{-2} \left( f(x) - 2 g(x) \right) dx =3$, find $\displaystyle\int_{-2}^{5} g(x) dx$
\end{document}